% The Golden-Ratio Polyhedral Loop: A Complete Forward Cycle with Algebraic and Geometric Closure
% Author: Shiva Meucci
% Version: v11 (February 2026)
%
% This is the source paper for the Lean 4 formalization in LeanProofing/.
% The LaTeX source is stored here for reference alongside the formal proofs.
%
% STRUCTURE (v11):
%   Section 1:  Reader's Guide (three levels of understanding)
%   Section 2:  Introduction (historical context, operations, journey overview,
%               main theorem preview, key discoveries)
%   Section 3:  Number-Theoretic Foundation
%               - Golden ratio phi = (1+sqrt5)/2, minimal polynomial x^2-x-1=0
%               - Lucas numbers L_n = phi^n + psi^n, recurrence, integrality
%               - Theorem 3.2: L_8 = phi^8 + phi^{-8} = 47 (THE critical identity)
%               - Proposition 3.1 (self-similar phase): {phi^n}/phi^{-n} = phi^n - 1
%               - cos(pi/5) = phi/2 (continuous-discrete bridge)
%   Section 4:  The Fundamental 3D Boundary
%               - Klein's theorem (1884): classification of finite rotation groups in SO(3)
%               - Corollary: I (icosahedral, order 60) is the largest
%               - Orbit-stabiliser theorem
%               - 60-face saturation proposition
%               - Why icosahedral symmetry is necessary for phi-geometry
%   Section 5:  The Algebraic Description of Closure
%               - Proposition 5.1: Each RT main-sequence stellation scales circumradius by phi^2
%               - Corollary 5.1: c_ell -> phi^{2*ell} * c_ell (degree-dependent amplification)
%               - Theorem 5.1: 4th stellation -> exactly 60 faces (geometric saturation)
%               - Remark 5.2: Why the main sequence stops at 4 (phi^8 < 60 < phi^10)
%               - Theorem 5.2: Algebraic closure -- phi^{8*ell} + phi^{-8*ell} = L_{8*ell} in Z
%               - Remark 5.4: Arithmetic signature -- 60 = L_8 + F_7 = 47 + 13
%               - Section 5.5: Octonionic mirror (Hurwitz 1898, E_8 lattice, 8-60 correspondence)
%   Section 6:  The Geometric Description of Closure
%               - Phase I  (Section 6.1): Platonic Spine
%                   T --(rect)--> O --(dual)--> C --(rect)--> CO --(dual)--> RD
%               - Phase II (Section 6.2): Transition to Icosahedral Symmetry
%                   T_emb --(I-orbit)--> 5T --(hull)--> D --(rect)--> ID --(dual)--> RT
%                   Key: Theorem 6.2 (Coset Completion -- T_emb -> 5T -> D)
%                   Key: Theorem 6.5 (Golden Rhombus -- RT faces have diagonal ratio phi:1)
%               - Phase III (Section 6.3): Golden Scaling to Maximum
%                   RT --(stell^4)--> RHC  (c_ell -> phi^{8*ell} * c_ell)
%               - Phase IV (Section 6.4): Forward Return
%                   RHC --(deform)--> DH --(extract)--> ID --(extract)--> T
%                   Key: Theorem 6.8 (Two Distinguished Members of hexecontahedral family)
%                   Key: Theorem 6.9 (DH Vertex Decomposition: 12+20+30=62)
%                   Key: Theorem 6.10 (Harmonic-Guided Extraction: DH->ID)
%                   Key: Theorem 6.11 (Compound of 5T in ID -> T extraction)
%               - Section 6.5: Symmetry Incommensurability and Quasicrystalline Closure
%                   Lemma 6.1: Pentagonal-tetrahedral incommensurability (5 not in T)
%                   Theorem 6.12: Topological closure without orientational periodicity
%   Section 7:  Spherical Harmonics and Poole's Invariants
%               - Poole (1932): invariant harmonics at degrees {0,6,10,12,15,...}
%               - P = Y_6^(I) (degree-6 generator), Q = Y_{10}^(I) (degree-10)
%               - LCM(6,10)=30 (faces of RT), 2*LCM(6,10)=60 (faces of RHC)
%               - Harmonic amplification under 4 stellations
%   Section 8:  The Belt: Geometric Closure and Quasicrystalline Structure
%               - Proposition 8.1: Belt width = 2*pi*phi^{-8} (in Wenninger coordinate model)
%               - Two interpretations: algebraic (fractional offset of phi^8) vs
%                 geometric (angular incommensurability)
%               - Connection to quasicrystalline tilings
%   Section 9:  Synthesis: The Complete Forward Loop
%               - Proof of Main Theorem (algebraic + geometric + convergence)
%               - Summary table of journey steps
%               - Theorem 9.1 (Distinguished Forward Loop): geometric necessity of each step
%               - Cascade of constraints (10 steps)
%               - Why these values cannot be otherwise
%   Section 10: Conclusion
%   Section 11: Further Directions (harmonic analysis and geometric topology)
%   Appendix A: Poole Invariant -- Symmetry-Based Critical-Point Analysis
%               - Explicit saddle-point verification via C_2 symmetry argument
%               - Hessian analysis: mixed signature at ID vertices
%               - Poincare-Hopf check: 12(+1) + 20(+1) + 30(-1) = 2 = chi(S^2)
%
% KEY CHANGES FROM v10b:
%   1. General Lucas closure theorem: phi^{8*ell}+phi^{-8*ell}=L_{8*ell} for all ell
%      (v10b only had the specific case ell=1, i.e., L_8=47)
%   2. Self-similar phase structure proposition (Section 3.4)
%   3. Degree-dependent harmonic amplification formalized as Corollary 5.1
%   4. Arithmetic remark: 60 = 47 + 13 = L_8 + F_7 (Section 5.4)
%   5. Cleaner separation of "algebraic" vs "geometric" descriptions throughout
%   6. Coset completion now called "coset activation" and treated as distinct operation
%   7. Poole invariant appendix is now more rigorous (Poincare-Hopf verification explicit)
%   8. Section numbers reorganised: Synthesis is now Section 9 (was implied Section 13 in v10b)
%   9. Octonionic section now clearly labelled as structural remark, not new theorem
%  10. Forward return steps now numbered Steps 7-9 (was 7-9 but unnumbered in some versions)
%
% REFERENCES:
%   Coxeter 1973, Wenninger 1971, Coxeter-Longuet-Higgins-Miller 1954,
%   Klein 1884, Poole 1932, Conway-Burgiel-Goodman-Strauss 2008,
%   Hurwitz 1898, Baez 2002, Elser-Sloane 1987, Koca-Koca-Koc 2011,
%   Kepler 1619, Shechtman et al. 1984, de Bruijn 1981,
%   Senechal 1995, Baake-Grimm 2013, Humphreys 1990
